\pdfoutput=1
\documentclass[11pt]{article}
\usepackage[T1]{fontenc}
\usepackage{lmodern}
\usepackage{amsmath,amssymb,amsthm,mathtools}
\usepackage{booktabs,array}
% Generated by scripts/make_claims_tex.py from claims.json -- do not edit by hand.
\newcommand{\clDate}{November 20, 2026}
\newcommand{\claSixteen}{37}
\newcommand{\claSeventeen}{42}
\newcommand{\claEighteen}{47}
\newcommand{\claNineteen}{52}
\newcommand{\claTwenty}{58}
\newcommand{\clEnumEleven}{696}
\newcommand{\clEnumTwelve}{4}
\newcommand{\clEnumThirteen}{2155}
\newcommand{\clEnumFourteen}{579}
\newcommand{\clEnumFifteen}{2298}
\newcommand{\clTorusTwenty}{51}
\newcommand{\clTorusTwentyT}{52}
\newcommand{\clTorusPairs}{163{,}664{,}896}
\newcommand{\clTorusSearched}{4675}
\newcommand{\clCubes}{4708}
\newcommand{\clSelftestTrials}{18{,}000}
\newcommand{\clSatSamples}{411}
\newcommand{\clBruteN}{8}
\newcommand{\clCpsatN}{7}
\newcommand{\clLpUniform}{2000}
\newcommand{\clLpBalanced}{2000}
\newcommand{\clLpNbSixteen}{529}
\newcommand{\clLpNbNineteen}{742}
\newcommand{\clSatN}{8}
\newcommand{\clNosymN}{11}
\newcommand{\clSatUniform}{200}
\newcommand{\clSatBalanced}{200}
\newcommand{\clSatCubesTwelve}{8}
\newcommand{\clSatCubesThirteen}{3}
\newcommand{\clTorusMinutes}{33}
\newcommand{\clWorkersMax}{8}
\newcommand{\clEnumSixteen}{26}
\newcommand{\clEnumLabSixteen}{416}
\newcommand{\clEnumCubesSixteen}{65{,}536}
\newcommand{\clEnumSeventeen}{5}
\newcommand{\clEnumLabSeventeen}{80}
\newcommand{\clSatSixteenLeaves}{400}
\newcommand{\clReplayNineteen}{24}
\newcommand{\clReplayTwenty}{24}
\newcommand{\clpdfEpoch}{1{,}795{,}132{,}800}
\newcommand{\clLpNbTwenty}{821}
\newcommand{\clPqvHoursNineteen}{290}
\newcommand{\clPqvHoursTwenty}{1400}
\newcommand{\clTorusTwentyTwo}{64}
\newcommand{\clTorusTwentyTwoT}{65}
\newcommand{\clTorusPairsTwentyTwo}{2{,}101{,}660{,}623}
\newcommand{\clTorusSearchedTwentyTwo}{80}
\newcommand{\clTorusTwentyTwoB}{66}
\newcommand{\clpages}{11}
\newcommand{\clLeavesSixteen}{204{,}573{,}866}
\newcommand{\clLeavesSixteenSci}{2.0\times10^{8}}
\newcommand{\clFlowsSixteenSci}{3.0\times10^{8}}
\newcommand{\clCpuSecSixteen}{1457}
\newcommand{\clCpuHSixteen}{0.4}
\newcommand{\clLeavesSeventeen}{817{,}258{,}210}
\newcommand{\clLeavesSeventeenSci}{8.2\times10^{8}}
\newcommand{\clFlowsSeventeenSci}{1.2\times10^{9}}
\newcommand{\clCpuSecSeventeen}{6489}
\newcommand{\clCpuHSeventeen}{1.8}
\newcommand{\clLeavesEighteen}{3{,}562{,}197{,}100}
\newcommand{\clLeavesEighteenSci}{3.6\times10^{9}}
\newcommand{\clFlowsEighteenSci}{5.6\times10^{9}}
\newcommand{\clCpuSecEighteen}{34{,}763}
\newcommand{\clCpuHEighteen}{9.7}
\newcommand{\clLeavesNineteen}{14{,}324{,}893{,}730}
\newcommand{\clLeavesNineteenSci}{1.4\times10^{10}}
\newcommand{\clFlowsNineteenSci}{2.3\times10^{10}}
\newcommand{\clCpuSecNineteen}{181{,}508}
\newcommand{\clCpuHNineteen}{50.4}
\newcommand{\clLeavesTwenty}{57{,}257{,}809{,}825}
\newcommand{\clLeavesTwentySci}{5.7\times10^{10}}
\newcommand{\clFlowsTwentySci}{9.4\times10^{10}}
\newcommand{\clCpuSecTwenty}{798{,}804}
\newcommand{\clCpuHTwenty}{221.9}
\newcommand{\clPqvPartsEighteen}{199}
\newcommand{\clPqvLeavesEighteen}{3{,}562{,}197{,}100}
\newcommand{\clPqvFracEighteen}{100.0}
\newcommand{\clLpMaxSixteen}{37.50}
\newcommand{\clLpMaxSeventeen}{42.11}
\newcommand{\clLpMaxEighteen}{47.50}
\newcommand{\clLpMaxNineteen}{52.56}
\newcommand{\clLpMaxTwenty}{58.44}
\newcommand{\clTSixteen}{38}
\newcommand{\clTSeventeen}{43}
\newcommand{\clTEighteen}{48}
\newcommand{\clTNineteen}{53}
\newcommand{\clTTwenty}{59}
\newcommand{\clGrowthMin}{4.0}
\newcommand{\clGrowthMax}{4.4}
\newcommand{\clRunRows}{3 & 1 & 2 & 3 & 0 & $<0.01$\,s\\
4 & 2 & 3 & 10 & 0 & $<0.01$\,s\\
5 & 4 & 5 & 21 & 0 & $<0.01$\,s\\
6 & 5 & 6 & 169 & 0 & $<0.01$\,s\\
7 & 7 & 8 & 686 & 0 & $<0.01$\,s\\
8 & 9 & 10 & 2\,450 & 0 & $0.01$\,s\\
9 & 12 & 13 & 8\,526 & 0 & $0.02$\,s\\
10 & 14 & 15 & 46\,481 & 0 & $0.22$\,s\\
11 & 17 & 18 & 189\,337 & 0 & $0.93$\,s\\
12 & 21 & 22 & 688\,447 & 0 & $2.34$\,s\\
13 & 24 & 25 & 2\,994\,131 & 0 & $13.0$\,s\\
14 & 28 & 29 & 12\,441\,976 & 0 & $58.2$\,s\\
15 & 32 & 33 & 51\,084\,014 & 0 & $303$\,s\\
16 & 37 & 38 & 204\,573\,866 & 0 & $1457$\,s\\
17 & 42 & 43 & 817\,258\,210 & 0 & $6489$\,s\\
18 & 47 & 48 & 3\,562\,197\,100 & 0 & $9.7$\,h\\
19 & 52 & 53 & 14\,324\,893\,730 & 0 & $50.4$\,h\\
20 & 58 & 59 & 57\,257\,809\,825 & 0 & $221.9$\,h\\}

\usepackage[margin=1in]{geometry}
\usepackage[colorlinks=true,urlcolor=blue,citecolor=blue,linkcolor=blue]{hyperref}

\newtheorem{theorem}{Theorem}[section]
\newtheorem{proposition}[theorem]{Proposition}
\newtheorem{lemma}[theorem]{Lemma}
\newtheorem{corollary}[theorem]{Corollary}
\theoremstyle{remark}
\newtheorem{remark}[theorem]{Remark}

\emergencystretch=3em
\newcommand{\W}{\mathsf W}
\newcommand{\B}{\mathsf B}

\title{Peaceable queens on boards of size 16 to 20}
\author{Jason Davies\thanks{\href{mailto:jason@jasondavies.com}{\nolinkurl{jason@jasondavies.com}}}}
\date{\clDate}

\begin{document}

\maketitle
\begin{abstract}
Let $a(n)$ be the largest $m$ such that $m$ white and $m$ black queens can be placed on an $n\times n$
board with no queen attacking a queen of the other colour (OEIS A250000). The OEIS gave $a(n)$ for
$n\le15$; DiMichele and Forbes determined $a(16)$ by integer programming in 2025, and in 2026 Arevalo
determined $a(17)$ and $a(18)$ in an unrefereed, AI-assisted computational note. We prove $a(19)=\claNineteen$ and $a(20)=\claTwenty$, which
were open, and we confirm $a(16)=\claSixteen$, $a(17)=\claSeventeen$ and $a(18)=\claEighteen$
independently, with two separately written programs. In each case $a(n)=\lfloor7n^2/48\rfloor$, so the
constructions that Ainley gave in 1977 are optimal on these boards. The proof is computer-assisted. As in
earlier work, we label each of the $6n-2$ lines of the board white or black, which turns the problem into
maximising the smaller of two cell counts. Our new ingredient is the bound: for fixed rows and columns we
bound that quantity by Lagrangian duality, and each Lagrangian subproblem is a minimum cut on the $4n-2$
diagonals. With symmetry breaking, the search for $n=20$ meets $\clLeavesTwentySci$ row/column labellings
up to symmetry, and a single family of minimum cuts prunes every one of them. That search took
$\clCpuHTwenty$ CPU hours, mostly on cloud machines. A second implementation repeats the whole search for
every $n\le18$. Orbit counts by Burnside's lemma confirm that every search is complete, and exact LP
certificates and SAT proofs checked by \texttt{drat-trim} give further checks. As side results we find
$\clEnumSixteen$ and $\clEnumSeventeen$ inequivalent optimal configurations for $n=16$ and $n=17$ (OEIS
A260680), and we confirm by computation that the $20\times20$ and $22\times22$ tori admit
$\clTorusTwenty$ and $\clTorusTwentyTwo$ queens of each colour but not more (A279405), values that also
follow from Harries's exact formula for even tori.
\end{abstract}

\noindent\textbf{AI transparency statement.} The results, the computer-assisted proofs, the code and this
manuscript were produced by artificial intelligence. The human author initiated the broader research and
directed the AI systems, but did not write the text and has not checked the arguments. The AI system's
literature searches missed the work of DiMichele and Forbes, of Arevalo and of Harries discussed below,
and earlier versions of this paper, and our edits of OEIS A250000, gave incorrect credit for
$a(16)$--$a(18)$. Details are given at the end of the paper.

\section{Introduction}

Two queens of different colours \emph{attack} each other if they share a row, a column or a diagonal and
no queen stands between them. The \emph{peaceable queens} problem asks for $a(n)$, the largest $m$ such
that $m$ white and $m$ black queens can be placed on an $n\times n$ board with no white queen attacking a
black queen. It is Problem~C1 of Ainley's \emph{Mathematical Puzzles}~\cite{Ainley}. Bosch posed it again
in 1999 and solved small cases by integer programming~\cite{Bosch1,Bosch2}, and it is sequence A250000 of
the OEIS~\cite{OEIS}.

Ainley gave constructions for every $n\le30$ and a general construction of asymptotic density $7/48$.
Ainley's values agree with $\lfloor7n^2/48\rfloor$ for $14\le n\le30$ except at $n=27$, where a board
with $106$ queens of each colour is now known; for every other $n$ with $16\le n\le30$ Ainley's values are
still the best known~\cite{OEIS}. The exact values
$a(1),\dots,a(15)=0,0,1,2,4,5,7,9,12,14,17,21,24,28,32$ were obtained with several methods. Smith, Petrie
and Gent~\cite{SPG} used constraint programming with symmetry breaking (see also~\cite{Petrie}), and
Prestwich and Beck~\cite{PB} obtained $a(12)$ and $a(13)$. In 2014 Pratt confirmed $a(1),\dots,a(13)$ by integer
programming and gave upper bounds such as $a(16)\le64$ and $a(19)\le90$. Tabatabai found $a(14)$ and
$a(15)$ by integer programming in 2018 (all three in~\cite{OEIS}). DiMichele and Forbes~\cite{DF}
determined $a(16)=37$ with an integer program over line variables, solved by Gurobi in about $40$ days
(published online in August 2025). The problem is also an exercise in
Knuth's volume on satisfiability~\cite{Knuth}.

Since $\lfloor 7n^2/48\rfloor$ matches every known value except $n=5$ and $n=9$ (where $a(n)$ is one
larger), it has been asked whether Ainley's construction is optimal. Yao and Zeilberger~\cite{YZ}
studied the problem numerically and symbolically, and Clinch, Drescher, Huynh and Saffidine~\cite{CDHS}
record their conjecture that it is. Clinch et al.\ also proved $a(n)\le0.1716\,n^2$ for large $n$,
improving the bound $n^2/4$ of van Bommel and MacEachern~\cite{vBM}, and they noted that only the first
fifteen terms were known. Kane
studied the asymptotics~\cite{Kane}.

\paragraph{Prior and concurrent work.} Describing a placement by labels on lines is not new. Smith,
Petrie and Gent~\cite[\S5]{SPG} used binary line variables (a line has no white queen, or no black
queen) and remarked that one could search on the $6n$ line variables alone. Prestwich and Beck~\cite{PB}
wrote an integer program with $0/1$ line variables, as does the tighter formulation of DiMichele and
Forbes~\cite{DF}, and the upper bounds of Clinch et al.~\cite{CDHS}
use the regions cut out by the lines that meet one army. Arevalo~\cite{Arevalo} states the line-labelling
reformulation of Lemma~\ref{lem:label} below, with a proof (Lemma~1 of~\cite{Arevalo}). With a search that also
takes rows and columns before diagonals, the product bound $m\le|R||C|$, a bound for each diagonal family
on its own, and the group $D_4\times\{\text{colour swap}\}$, Arevalo determined $a(17)=42$ and $a(18)=47$, and
independently $a(16)=37$, in August and September 2026. The note is AI-assisted and unrefereed; its runs for $n=17$ and $18$ use a
single engine, and its DRUP proofs go up to $n=7$. It lists $a(19)$ as future work, and it asks, as an open
question, for a joint bound over both diagonal families. Our minimum-cut Lagrangian bound
(Lemmas~\ref{lem:lag} and~\ref{lem:cut}) is such a joint bound; we have not found it elsewhere. On the
torus, Harries~\cite{Harries} proved an exact formula $t(2q)=H(2q)$ for even orders, with a Lean proof of
the main theorem, using a line-colouring equivalence (Lemma~6 of~\cite{Harries}) and a partially dualized
Lagrangian. The formula gives the torus values for $n=20$ and $22$ in Section~\ref{sec:related}, and
Harries also found
$t(15)=20$ and $t(17)=28$. Thus our new results are $a(19)$ and $a(20)$, the counts A260680$(11..17)$ and
the minimum-cut bound; for $a(16)$--$a(18)$ and the torus our computations are independent confirmations.
Earlier versions of this paper, and our edits of OEIS A250000, first presented $a(16)$--$a(18)$ as new and
then credited $a(16)$ to Arevalo, because the AI system's literature searches missed these works.

\begin{theorem}\label{thm:main}
$a(16)=\claSixteen$, $a(17)=\claSeventeen$, $a(18)=\claEighteen$, $a(19)=\claNineteen$ and
$a(20)=\claTwenty$.
\end{theorem}

The lower bounds come from known constructions. For $n=16$ we use the Ainley arrangement shown in the
OEIS entry (Figure~\ref{fig:16}). For $n=17,\dots,20$ we use the boards in Kamenetsky's file of best known
solutions~\cite{OEIS}. We checked each of them directly for attacks. The new content of
Theorem~\ref{thm:main} is the upper bounds for $n=19$ and $20$; for $n=16,17,18$ the upper bounds confirm
those of DiMichele and Forbes~\cite{DF} ($n=16$) and Arevalo~\cite{Arevalo} ($n=16,17,18$). The upper bounds say that no placement has $\clTSixteen$, $\clTSeventeen$, $\clTEighteen$, $\clTNineteen$ or
$\clTTwenty$ queens of each colour, respectively.

\begin{figure}[htbp]
\centering
{\small\ttfamily
\begin{tabular}{l}
. . . . W W W W . . . . . . . .\\
. . . . W W W W . . . . . . . W\\
. . . . W W W W . . . . . . W W\\
. . . . W W W W . . . . . W W W\\
. . . . . W W . . . . . W W W W\\
. . . . . . . . . . . . W W W W\\
. . . . . . . . . . . . W W W .\\
. . . . . . . . . . . . W W . .\\
. . . B . . . . . . . . . . . .\\
. . B B . . . . . . . . . . . .\\
. B B B . . . . . . . . . . . .\\
B B B B . . . . . . B B . . . .\\
B B B B . . . . . B B B . . . .\\
B B B . . . . . B B B B . . . .\\
B B . . . . . . B B B B . . . .\\
B . . . . . . . B B B B . . . .
\end{tabular}}
\caption{An Ainley-type arrangement of $37$ white and $37$ black queens on the $16\times16$ board, as
reproduced in~\cite{OEIS}.}\label{fig:16}
\end{figure}

\section{Line labellings}

The board has $n$ rows, $n$ columns, $2n-1$ diagonals ($i-j$ constant) and $2n-1$ antidiagonals
($i+j$ constant), so $6n-2$ \emph{lines} in all. Every cell lies on exactly four lines, and distinct
cells have distinct (diagonal, antidiagonal) pairs. A \emph{labelling} $x$ assigns $\W$ or $\B$ to each
line. A cell is a \emph{$\W$-cell} of $x$ if all four of its lines are labelled $\W$, and a
\emph{$\B$-cell} if all four are labelled $\B$. Let $W(x)$ and $B(x)$ be the numbers of $\W$-cells and
$\B$-cells. No cell is both.

\begin{lemma}\label{lem:label}
$a(n)=\max_x\min\{W(x),B(x)\}$.
\end{lemma}

\begin{proof}
Take a peaceable placement with $m$ queens of each colour. Suppose some line contains queens of both
colours. Walking along it, we meet two consecutive queens of different colours, with no queen between
them, and these two attack each other. So every line contains queens of at most one colour. Label a line
by the colour of its queens, and label empty lines arbitrarily. Then every white queen stands on a
$\W$-cell and every black queen on a $\B$-cell, so $\min\{W,B\}\ge m$. Conversely, given $x$, put
$m=\min\{W,B\}$ white queens on $\W$-cells and $m$ black queens on $\B$-cells. A line containing queens of
both colours would carry both labels, which is impossible. So no white queen shares a line with a black
queen, and in particular none attacks one.
\end{proof}

This reformulation is known: it is Lemma~1 of Arevalo~\cite{Arevalo}, its torus form is Lemma~6 of
Harries~\cite{Harries}, and binary line variables go back to~\cite{SPG,PB} (see also~\cite{DF}). The proof shows that it does
not matter whether ``attack'' allows blocking. So the question is whether
some labelling has $W\ge T$ and $B\ge T$, where $T=\lfloor7n^2/48\rfloor+1$.

\section{The bound}

Let $R$ and $C$ be the sets of rows and columns labelled $\W$. Every $\W$-cell lies in $R\times C$ and
every $\B$-cell in $R^c\times C^c$. This gives the \emph{count bound}: if some rows and columns are
labelled, then $W\le r_{\W} c_{\W}$, where $r_{\W}$ and $c_{\W}$ are the numbers of rows and columns not
labelled $\B$. The same bound with the colours exchanged holds for $B$. Once $(R,C)$ is fixed, call the
cells of $R\times C$ \emph{$\W$-candidates} and those of $R^c\times C^c$ \emph{$\B$-candidates}. Write
$x_d,y_e\in\{0,1\}$ for the labels of diagonal $d$ and antidiagonal $e$, with $1$ meaning $\W$. A
$\W$-candidate on $d$ and $e$ is a $\W$-cell if and only if $x_d=y_e=1$. A $\B$-candidate is a $\B$-cell
if and only if $x_d=y_e=0$.

\begin{lemma}[Lagrangian bound]\label{lem:lag}
For any set $X$ of labellings and any $\lambda\in[0,1]$,
\[\max_{x\in X}\min\{W(x),B(x)\}\le \max_{x\in X}\{\lambda W(x)+(1-\lambda)B(x)\}.\]
\end{lemma}

\begin{proof}
$\min\{W,B\}\le\lambda W+(1-\lambda)B$ holds for each $x$.
\end{proof}

\begin{lemma}[Min-cut evaluation]\label{lem:cut}
Fix $(R,C)$, integers $0\le p\le Q$ and any partial labelling of the diagonals and antidiagonals. There
is a network $N_p$, on a source $s$, a sink $t$ and one node for each unlabelled diagonal or
antidiagonal, and a constant $G$, such that each completion $x$ satisfies
$pW(x)+(Q-p)B(x)=G-\operatorname{cap}(S_x)$, where $S_x$ is the $s$--$t$ cut whose source side
consists of $s$ and the lines labelled $\W$ by $x$. Hence
$U_p:=\max_x\{pW+(Q-p)B\}=G-\operatorname{mincut}(N_p)$.
\end{lemma}

\begin{proof}
For $0/1$ variables,
\begin{align*}
p\,xy&=p-p\,[x=0]-p\,[x=1,y=0],\\
(Q-p)(1-x)(1-y)&=(Q-p)-(Q-p)\,[x=1]-(Q-p)\,[x=0,y=1].
\end{align*}
With the source side meaning $\W$, the four bracketed indicators are exactly the costs of the arcs
$s\to x$, $x\to y$, $x\to t$ and $y\to x$, each of capacity $p$ or $Q-p$, in the cut $S_x$. We sum these
identities over the $\W$-candidates (with $x=x_d$, $y=y_e$) and the $\B$-candidates. If a line is
already labelled, we substitute its value. The term then becomes a constant (its gain, or $0$) or a
single arc at the other line. Every node has arcs from $s$ and to $t$, and exactly one of these two is
cut. So we may subtract the smaller of the two capacities from both and from $G$. Distinct completions
give distinct cuts, and every cut arises from a completion, so the maximum over completions is
$G$ minus the minimum cut.
\end{proof}

This bound couples the two diagonal families, which a bound for each family separately does not
(compare~\cite[\S7]{Arevalo}); we have not found it in the literature. Equivalently, $pW+(Q-p)B$ is a
supermodular pseudo-Boolean function of the free lines in which every
term involves at most two variables. Such functions are maximised by minimum cuts
\cite{Hammer,KZ}. Without the row/column stage every cell would couple four lines, and the same
function would be a maximum-weight closure~\cite{Picard}. The program uses $Q=256$.

\begin{corollary}[soundness needs only a feasible flow]\label{cor:sound}
If $f$ is the value of any feasible $s$--$t$ flow in $N_p$ and $G-f<QT$, then no completion has
$W\ge T$ and $B\ge T$.
\end{corollary}

\begin{proof}
Any feasible flow has value at most the capacity of every cut, so for every completion $x$ we have
$pW+(Q-p)B=G-\operatorname{cap}(S_x)\le G-f<QT$. By Lemma~\ref{lem:lag} with $\lambda=p/Q$,
$\min\{W,B\}<T$.
\end{proof}

So the correctness of the pruning does not depend on the max-flow routine finding an optimal flow or on
how $p$ is chosen. These only affect how much is pruned. The program picks $p$ by bisection: $U_p$ is
convex in $p$, and $W(x^*)-B(x^*)$ is a subgradient at the optimal completion $x^*$. The completion
$x^*$ is itself a labelling, and if $\min\{W(x^*),B(x^*)\}\ge T$ the search has found a solution.

\section{Symmetry and search}

The group $G=D_4\times\{\text{colour swap}\}$ has order~$16$ and acts on labellings. A symmetry of the
board maps the four lines of each cell to the four lines of the image cell, and the colour swap
exchanges $W$ and $B$. So $\min\{W,B\}$ and the conditions $W,B\ge T$ are $G$-invariant. The row and
column lines are enumerated in groups $\{\text{row }e,\text{ row }n-1-e,\text{ column }e,\text{ column
}n-1-e\}$ for $e=0,1,\dots$; when $n$ is odd the last group is $\{\text{row }m,\text{ column }m\}$ with
$m=(n-1)/2$. Each group is mapped onto itself by every element of $G$.

\begin{lemma}\label{lem:sym}
Order the row and column lines group by group. It suffices to search the labellings whose row/column part
is lexicographically minimal in its $G$-orbit. Moreover, at the end of every group the search may
discard any prefix that is not lexicographically $\le$ each of its $16$ images.
\end{lemma}

\begin{proof}
Let $S$ be the set of row and column lines. It is $G$-invariant, so $G$ acts on labellings of $S$. Each
orbit of labellings contains one whose restriction to $S$ is lexicographically minimal among the
restrictions of its images, and that labelling has the same value of $\min\{W,B\}$. Let $P$ be a union of
initial groups. It is $G$-invariant, so the restriction of $g\cdot x$ to $P$ depends only on the
restriction of $x$ to $P$. The order lists $P$ first, so lexicographic comparison of whole sequences
compares their restrictions to $P$ first. If $x|_S\le(g\cdot x)|_S$ then $x|_P\le(g\cdot x)|_P$. So
the lex-minimal representative survives every prefix test.
\end{proof}

The test at the end of the last group keeps a labelling of $S$ only if it is the lex-minimum of its
orbit, so at most one labelling per orbit survives. Lemma~\ref{lem:sym} says that exactly one does. We
also check this numerically. By Burnside's lemma we count the orbits of pairs $(R,C)$ that pass the
count bound, and this number must equal the number of leaves of the search (Section~\ref{sec:ver}).

The search has two stages. Stage~1 is a depth-first search over row and column labels with the pruning
of Lemma~\ref{lem:sym} and the count bound. At each complete $(R,C)$, Stage~2 evaluates
Lemma~\ref{lem:cut} with all diagonals free. If no $p$ prunes, it branches on the diagonals and
antidiagonals and recomputes the bound at every node. For $n\ge16$ the work is split into \emph{cubes}:
the canonical prefixes of the first four groups (sixteen lines). There are $\clCubes$ of them for each
$n=16,\dots,20$. Each cube is solved by a separate process and its result is logged together with its
node and flow counts.

\section{Computation}

\begin{table}[!ht]
\centering
\begin{tabular}{rrrrrr}
\toprule
$n$ & $a(n)$ & $T$ & $(R,C)$ orbits & diagonal searches & CPU time\\\midrule
\clRunRows
\bottomrule
\end{tabular}
\caption{Searches proving that no labelling has $W,B\ge T$, with $T=a(n)+1$. The fourth column counts the
$(R,C)$ orbits that pass the count bound; each was pruned by Corollary~\ref{cor:sound} at the root, so
the fifth column, which counts the orbits whose diagonals had to be searched, is $0$ throughout. Times are
total CPU times, on a shared local machine (at most $\clWorkersMax$ processes at a time) and, for
$n=19,20$, mostly on cloud containers. For $n\le15$
the same program also finds a placement with $a(n)$ queens of each colour.}
\label{tab:runs}
\end{table}

Table~\ref{tab:runs} lists the runs. For $n=20$ and $T=\clTTwenty$ the $\clCubes$ cubes contain
$\clLeavesTwenty$ orbits of $(R,C)$ that pass the count bound. The root bound pruned all of them,
using $\clFlowsTwentySci$ maximum-flow computations, so no diagonal search was needed. For $n=19$ and
$n=20$ most cubes were run on rented cloud containers (Modal), with the same source code compiled there
by gcc~12.2; there a cube takes about $1.3$ times as long as on the local machine. Every line returned
from the cloud was checked (cube number, prefix and format) before it was logged, and random cubes were
rerun locally with identical results. From $n=15$ on, the number of orbits grows by a factor between
$\clGrowthMin$ and $\clGrowthMax$ per unit increase of $n$, and the running time roughly in proportion. The same program reproduces every known value $a(n)$, $n\le15$: it finds a placement
with $a(n)$ queens of each colour, and it proves that $a(n)+1$ is impossible.

\section{Verification}\label{sec:ver}

We checked the computation in the following ways. The first two use different code but the same method;
the others avoid parts of the method.

\paragraph{A second implementation.} It is written in Rust, separately from the first program (in C).
It differs in four ways:
\begin{itemize}
\item it enumerates pairs of bitmasks $(R,C)$ and keeps a pair if its integer code is minimal among its
$16$ images, a different canonical form;
\item it evaluates the Lagrangian bound on the generic closure network, with one auxiliary node per
candidate cell;
\item it solves that network with FIFO push--relabel~\cite{GT}, with $\lambda=p/1000$;
\item it searches any unpruned $(R,C)$ by plain depth-first search that uses only counting bounds.
\end{itemize}
Its pruning is sound by the analogue of Corollary~\ref{cor:sound}. For $n=16$ and $n=17$ it repeated the
whole search. It reported infeasibility, and its orbit counts ($\clLeavesSixteen$ and
$\clLeavesSeventeen$) agree exactly with those of the first program. For $n=18$ we ran it on slices:
slice $j$ of $m=199$ consists of the canonical pairs with $R\equiv j\pmod{199}$, where bit $i$ of $R$ is
row $i$, and the slices $j=0,\dots,198$ partition the search. All $\clPqvPartsEighteen$ slices were run
and every orbit was pruned; together they contain $\clPqvLeavesEighteen$ orbits, exactly the orbit count of
the first program and of Burnside's lemma. So for $n\le18$ two independent implementations cover the whole
search. The second implementation was not run for $n=19$ and $n=20$; we estimate that it would need about
$\clPqvHoursNineteen$ and $\clPqvHoursTwenty$ CPU hours on our machine.

\paragraph{Orbit counts by Burnside's lemma.} A short independent script counts, for every $n\le20$, the
$G$-orbits of pairs $(R,C)$ with $|R||C|\ge T$ and $|R^c||C^c|\ge T$. It uses Burnside's lemma and the
action of $G$ on pairs (reversals, transposition and complementation). The script checks the action
against the cell maps and checks the orbit counts against brute-force enumeration for $n\le8$. For
$n=16,\dots,20$ the counts are $\clLeavesSixteen$, $\clLeavesSeventeen$, $\clLeavesEighteen$,
$\clLeavesNineteen$ and $\clLeavesTwenty$. They agree with the total numbers of leaves over all cubes of the first program, and
for $n\le15$ with the first program and the second implementation. So the cubes together visit exactly one
pair from each orbit, and none is missing.

\paragraph{Exact LP certificates.} For fixed $(R,C)$, the LP relaxation
\begin{align*}
\max\ s\quad\text{s.t.}\quad &s\le\textstyle\sum_{\W\text{-cand}}z_c,\quad s\le\sum_{\B\text{-cand}}u_c,\\
&z_c\le x_d,\ z_c\le y_e\quad\text{for each $\W$-candidate $c$},\\
&u_c\le1-x_d,\ u_c\le1-y_e\quad\text{for each $\B$-candidate $c$},
\end{align*}
with $x_d,y_e,z_c,u_c\in[0,1]$ and $s\ge0$,
bounds $\max\min\{W,B\}$ over the diagonal labels. This uses neither supermodularity nor minimum cuts.
We solve it in floating point with HiGHS~\cite{HiGHS} through SciPy~\cite{SciPy}, round the dual
multipliers to rationals and repair them to a feasible dual solution. Its objective, evaluated in exact
rational arithmetic, is then a valid upper bound by weak duality. For each $n=16,\dots,20$ we certified
$\clLpUniform$ uniformly random pairs $(R,C)$ passing the count bound and $\clLpBalanced$ random
``balanced'' pairs (both products $\ge T+n$). We also certified the pair $(R,C)$ of the best known board
(rows and columns holding a white queen are labelled $\W$) and all pairs within Hamming distance~$2$ of
it, from $\clLpNbSixteen$ pairs for $n=16$ to $\clLpNbTwenty$ for $n=20$. Every
certified bound was below $T$. The largest bounds for $n=16,\dots,20$ were
$\clLpMaxSixteen$, $\clLpMaxSeventeen$, $\clLpMaxEighteen$, $\clLpMaxNineteen$ and $\clLpMaxTwenty$,
against $T=\clTSixteen,\clTSeventeen,\clTEighteen,\clTNineteen,\clTTwenty$.

\paragraph{SAT with checked proofs.} We used two encodings, each with sequential-counter cardinality
constraints, and the Glucose solver through PySAT~\cite{PySAT}. The first encodes placements directly by
cell variables, with a conflict clause for every pair of cells sharing a line, so it does not use
Lemma~\ref{lem:label}. The second is the line-labelling encoding. Every unsatisfiability claim was checked
with \texttt{drat-trim}~\cite{DRAT}. Both encodings confirm $a(n)$ for $n\le\clSatN$. In addition,
random canonical $(R,C)$ for $n=16$ were checked by SAT: $\clSatUniform$ uniform and $\clSatBalanced$
balanced ones (both products at least $56$). We also checked $\clSatCubesTwelve$ random cubes for $n=12$,
$T=22$, and $\clSatCubesThirteen$ for $n=13$, $T=25$, in which the first $12$ row and column labels are
fixed. All were unsatisfiable with verified proofs. These checks
do not depend on Lemmas~\ref{lem:lag}--\ref{lem:sym}.

\paragraph{The definition itself.} A depth-first search over sets of white queens finds $a(n)$ for
$n\le\clBruteN$ directly from the definition, using only $f(S)\ge T$, where $f(S)$ is the number of
cells sharing no line with the white set~$S$. A search over all $3^{n^2}$ placements with the blocking
rule of chess gives the same values for $n\le4$. A direct CP-SAT model~\cite{ORTools} of placements gives
the same values for $n\le\clCpsatN$.

\paragraph{Unit tests and replays.} The min-cut evaluation of Lemma~\ref{lem:cut} agrees with brute-force
maximisation over all completions on $\clSelftestTrials$ random instances with $3\le n\le8$, including partially
labelled diagonals. The same holds for the second implementation's network on random instances with
$n\le6$. Switching off the symmetry breaking gives identical answers for $n\le\clNosymN$, and so does
switching off the min-cut bound in the second implementation (plain search). Random cubes of every run were rerun
and gave log lines identical to the recorded ones, apart from the time.

\paragraph{Placements.} Every placement that we report was checked directly for attacks.

\paragraph{Review.} An AI referee, run separately from the sessions that designed the method and wrote the
programs, read the proofs
of Lemmas~\ref{lem:label}--\ref{lem:sym} and both programs adversarially. It wrote the Burnside, LP,
brute-force and CP-SAT checks above. It found no error affecting the results. Its minor corrections (a
misstated number of fixed lines for the $n=13$ SAT cubes in an earlier draft, Corollary~\ref{cor:sound}
as the precise statement of what the pruning needs, and the Burnside test of completeness) are
incorporated. Three rounds of external AI review (ChatGPT Pro) of the code also reran the search for
$n=16$, recomputed the Burnside counts for $n\le20$ and for the torus runs, and recomputed the largest
LP certificates for $n=20$. The first two rounds found thirteen defects, now corrected, none of which
affected the claimed values; the third confirmed the corrections. The resumable drivers trusted logged results in a reused directory whenever
the cube list matched, even for a different $T$; every run reported here used a fresh directory for its
$(n,T)$, and the drivers now bind each directory to a manifest of all parameters and refuse a mismatch.
The torus program shifted signed integers beyond their width while building symmetry representatives
(undefined behaviour); with the corrected program the torus runs for $n=20$ and $n=22$ give identical
counters. The verification script ignored a failure of the leaf sampler and accepted misspelled stage
names; the completion log of the A260680(16) run was missing from the repository; the pinned
\texttt{drat-trim} did not compile under GCC~14; the LP above was stated with $s\in[0,1]$; and the
solvers accepted out-of-range arguments (a huge $Q$ overflowed, and an empty part of a partition reported
infeasibility). The second round found that two drivers starting at the same moment in one directory
could publish each other's manifest (creation is now locked and atomic), that $K=0$ produced an empty
list of cubes which a driver reported as infeasible (now rejected, and no driver reports a verdict for an
empty work list), that the cube-range argument accepted trailing characters, and that the torus checks
did not require each part exactly once; it also asked for a stale sentence in the README to be updated. Auditing for the same patterns, we also made
the second implementation reject a slice that cannot contain any row mask. Each correction has a regression test in \texttt{scripts/verify.py}.

\paragraph{What has to be trusted.} The upper bounds rest on Lemmas~\ref{lem:label}--\ref{lem:sym} and
on the first program. Table~\ref{tab:cover} lists what checked each value. For $n\le18$ a second,
separately written program repeats the whole search with the same result. For $n=19$ and $n=20$ only the first program covers the whole search. There the Burnside counts
confirm that it visits every orbit; a sample of the cubes computed in the cloud was rerun locally with
identical results; and the exact LP certificates and SAT proofs on samples check the bound without
Lemmas~\ref{lem:lag}--\ref{lem:sym}. A self-contained proof certificate would give a value of $p$ and a
feasible flow for each of the $\clLeavesTwentySci$ orbits for $n=20$. That would take hundreds of
gigabytes. Instead each cube is logged with its counts, and any cube can be rerun in seconds to minutes.

\begin{table}[!ht]
\centering\small
\begin{tabular}{l>{\raggedright}p{2.1cm}>{\raggedright}p{3.4cm}>{\raggedright\arraybackslash}p{7.2cm}}
\toprule
$n$ & first program & second implementation & other checks\\\midrule
$\le15$ & complete & complete & definition search ($n\le\clBruteN$), CP-SAT ($n\le\clCpsatN$), SAT+DRAT ($n\le\clSatN$)\\
16 & complete & complete & Burnside; SAT+DRAT on \clSatSixteenLeaves{} leaves; LP; replays\\
17 & complete & complete & Burnside; LP certificates; replays\\
18 & complete & complete (\clPqvPartsEighteen{} slices) & Burnside; LP certificates; replays\\
19 & complete & --- & Burnside; LP certificates; \clReplayNineteen{} cloud cubes rerun locally\\
20 & complete & --- & Burnside; LP certificates; \clReplayTwenty{} cloud cubes rerun locally\\
\bottomrule
\end{tabular}
\caption{What checked each upper bound. ``Burnside'' means that the number of orbits visited equals the
Burnside count; ``LP'' the exact LP certificates on sampled $(R,C)$; ``replays'' reruns of random cubes
with identical log lines.}
\label{tab:cover}
\end{table}

\section{Related sequences}\label{sec:related}

The same method gives side results on two related sequences. Each was computed by a single
implementation only, and each reproduces all the known terms of its sequence. The counts for A260680 are
new as far as we know (the OEIS gives them up to $n=10$); the torus values confirm known ones.
\begin{itemize}
\item \emph{Counting optimal configurations (A260680).} Two solutions are equivalent if they differ by an
element of $G$. We enumerate, without symmetry breaking, all labellings with $\min\{W,B\}\ge a(n)$. We
expand each into its placements and count orbits under $G$. This reproduces the known terms
$1,1,1,10,3,35,19,71,18,380$ for $n\le10$ and gives $\clEnumEleven$, $\clEnumTwelve$, $\clEnumThirteen$,
$\clEnumFourteen$ and $\clEnumFifteen$ for $n=11,\dots,15$. For $n=16$ the enumeration, split into
$\clEnumCubesSixteen$ cubes, finds $\clEnumLabSixteen$ labellings with $\min\{W,B\}\ge37$, all with
different cell sets, and $\clEnumSixteen$ inequivalent optimal configurations. Since
$\clEnumLabSixteen=16\cdot\clEnumSixteen$, no optimal configuration for $n=16$ is fixed by a nontrivial
element of~$G$. For $n=17$ the same enumeration finds $\clEnumLabSeventeen$ labellings with
$\min\{W,B\}\ge42$ and $\clEnumSeventeen$ inequivalent optimal configurations, again
$\clEnumLabSeventeen=16\cdot\clEnumSeventeen$. The split version reproduces the values for $n=10,12,13,14$
first.
\item \emph{The torus (A279405).} On the $n\times n$ torus the lines are the rows, the columns and the
diagonals $i\pm j\bmod n$. The same method, with the subgroup of symmetries generated by translations,
the reflections $i\mapsto-i$ and $j\mapsto-j$, transposition and the colour swap, reproduces the known
values for $4\le n\le14$, $n=16$ and $n=18$, each at $a(n)$ and $a(n)+1$. For $n=20$ it gives $\clTorusTwenty$, the value that Harries's formula $t(2q)=H(2q)$~\cite[Cor.~4]{Harries}
already gives: a placement with
$\clTorusTwenty$ of each colour exists, and $\clTorusTwentyT$ is impossible ($\clTorusPairs$ orbits, of
which $\clTorusSearched$ needed a diagonal search, in $\clTorusMinutes$ CPU minutes). The repository version of the program, which can split the search into
parts, caches the representatives of complements and uses unsigned masks, reproduces all these values and
all the $n=20$ counters exactly. For $n=22$ it gives $\clTorusTwentyTwo=H(22)$~\cite[Table~1]{Harries}: it finds a placement with
$\clTorusTwentyTwo$ white and $\clTorusTwentyTwoB$ black queens, which we checked directly for toroidal
attacks, and it shows that $\clTorusTwentyTwoT$ of each colour is impossible ($\clTorusPairsTwentyTwo$
orbits, of which $\clTorusSearchedTwentyTwo$ needed a diagonal search). These two computations are
therefore independent confirmations by a different method, not new values; Harries also determined
$t(15)=20$ and $t(17)=28$. The value for $n=19$ is still
unknown. Odd tori are harder for the method, because the root bound fails for many orbits.
\end{itemize}

\paragraph{Availability.} The code, logs and verification scripts are in the accompanying
repository\footnote{\url{https://github.com/jasondavies/peaceable-queens}.}. The script
\texttt{scripts/verify.py} rebuilds the programs and checks every number quoted here. By default this
takes a few minutes. Its \texttt{-{}-full} option reruns the complete search for any $n$. The cloud runs used
the dispatcher in \texttt{code/modal/}, whose returned results are logged in the same format as the local
runs.

\section*{Declaration of AI-generated research and authorship}
This manuscript is the product of an AI-driven research process. The human author initiated a broader
investigation of open problems in the OEIS, and AI systems worked on this problem within it. AI systems
produced all of the following: the method, the proofs, the code, the computations and their checks, and
the organisation and writing of the text. The system was Claude Opus~5.5 (Anthropic), used through
Claude Code. The human author directed the AI systems but wrote no part of the
manuscript; the human author also started some of the cloud computations, using commands prepared by
the AI system.
The AI system's literature searches missed the work of DiMichele and Forbes~\cite{DF}, of
Arevalo~\cite{Arevalo} and of Harries~\cite{Harries}, which later checks found; this version credits it,
and corrects the earlier versions' claims that $a(16)$--$a(18)$ and the torus values for $n=20,22$ were
new and that Arevalo first determined $a(16)$.

The computational claims are reproduced by \texttt{scripts/verify.py}. At the time of writing, the
mathematics and the code have been checked only by AI systems and by computer verification. These checks
include an adversarial review of the proofs and of both search programs by an independent AI referee,
who also wrote the Burnside, LP-certificate, brute-force and CP-SAT checks of Section~\ref{sec:ver}, and three
rounds of external AI review of the code (ChatGPT Pro), whose findings are corrected
(Section~\ref{sec:ver}). No
independent human expert has reviewed the work. The human author has not read this manuscript or checked
its arguments. The author takes responsibility for the accuracy of this disclosure.

\raggedright
\begin{thebibliography}{99}
\bibitem{Ainley} S.~Ainley, \emph{Mathematical Puzzles}, G.~Bell \& Sons, London, 1977.
\bibitem{Arevalo} F.~Arevalo, \emph{The sixteenth, seventeenth and eighteenth peaceable queens numbers: 37,
42, 47}, computational research note (AI-assisted), August~17, 2026, updated September~3 and~7, 2026,
\url{https://fabianarevalo.com/peaceable-queens}; code at
\url{https://github.com/farev/Matematica/tree/main/conjectures/peaceable-queens}.
\bibitem{Bosch1} R.~A.~Bosch, \emph{Peaceably coexisting armies of queens}, Optima (Mathematical
Programming Society Newsletter) \textbf{62} (1999).
\bibitem{Bosch2} R.~A.~Bosch, \emph{Armies of queens, revisited}, Optima (Mathematical Programming Society
Newsletter) \textbf{64} (2000), 15.
\bibitem{CDHS} K.~Clinch, M.~Drescher, T.~Huynh, A.~Saffidine, \emph{Constructions, bounds, and algorithms
for peaceable queens}, in: Proceedings of the Symposium on Algorithm Engineering and Experiments
(ALENEX 2025), SIAM, 2025, 160--167.
\href{https://doi.org/10.1137/1.9781611978339.13}{doi:10.1137/1.9781611978339.13}; arXiv:2406.06974,
\href{https://doi.org/10.48550/arXiv.2406.06974}{doi:10.48550/arXiv.2406.06974}.
\bibitem{DF} N.~DiMichele, M.~Forbes, \emph{Puzzle---Queen puzzles: an educational approach to integer
programming techniques}, INFORMS Trans.\ Educ.\ \textbf{26}(3) (2026), 247--253; published online
August~1, 2025. \href{https://doi.org/10.1287/ited.2024.0122}{doi:10.1287/ited.2024.0122}.
\bibitem{GT} A.~V.~Goldberg, R.~E.~Tarjan, \emph{A new approach to the maximum-flow problem}, J.\ ACM
\textbf{35} (1988), 921--940. \href{https://doi.org/10.1145/48014.61051}{doi:10.1145/48014.61051}.
\bibitem{Harries} J.~P.~Harries, \emph{Peaceable queens on even tori: an exact parity formula}, preprint,
September~6, 2026 (first version July~27, 2026), with a Lean proof of the main theorem,
\url{https://github.com/jphme/math-problems/tree/main/peaceable-queens-even-torus}.
\bibitem{Hammer} P.~L.~Ivănescu (Hammer), \emph{Some network flow problems solved with pseudo-Boolean
programming}, Oper.\ Res.\ \textbf{13} (1965), 388--399.
\href{https://doi.org/10.1287/opre.13.3.388}{doi:10.1287/opre.13.3.388}.
\bibitem{HiGHS} Q.~Huangfu, J.~A.~J.~Hall, \emph{Parallelizing the dual revised simplex method}, Math.\
Program.\ Comput.\ \textbf{10} (2018), 119--142.
\href{https://doi.org/10.1007/s12532-017-0130-5}{doi:10.1007/s12532-017-0130-5}.
\bibitem{PySAT} A.~Ignatiev, A.~Morgado, J.~Marques-Silva, \emph{PySAT: a Python toolkit for prototyping
with SAT oracles}, in: SAT 2018, Lecture Notes in Comput.\ Sci.\ \textbf{10929}, Springer, 2018,
428--437. \href{https://doi.org/10.1007/978-3-319-94144-8_26}{doi:10.1007/978-3-319-94144-8\_26}.
\bibitem{Kane} D.~M.~Kane, \emph{Asymptotic results for the queen packing problem}, preprint (2017),
arXiv:1703.04538. \href{https://doi.org/10.48550/arXiv.1703.04538}{doi:10.48550/arXiv.1703.04538}.
\bibitem{Knuth} D.~E.~Knuth, \emph{The Art of Computer Programming, Volume~4, Fascicle~6:
Satisfiability}, Addison-Wesley, 2015, p.~180, exercise~488, and pp.~282--283.
\bibitem{KZ} V.~Kolmogorov, R.~Zabih, \emph{What energy functions can be minimized via graph cuts?},
IEEE Trans.\ Pattern Anal.\ Mach.\ Intell.\ \textbf{26} (2004), 147--159.
\href{https://doi.org/10.1109/TPAMI.2004.1262177}{doi:10.1109/TPAMI.2004.1262177}.
\bibitem{OEIS} OEIS Foundation Inc., \emph{The On-Line Encyclopedia of Integer Sequences}, sequences
A250000 (revision~379, including the comments, the extensions, and D.~Kamenetsky's file of best known
solutions for $12\le n\le30$), A260680 and A279405, \url{https://oeis.org}.
\bibitem{ORTools} L.~Perron, F.~Didier, \emph{CP-SAT}, OR-Tools version~9.15, Google,
\url{https://developers.google.com/optimization}.
\bibitem{Petrie} K.~E.~Petrie, \emph{Symmetry breaking in peaceably coexisting armies of queens}, in: CP 2002,
Lecture Notes in Comput.\ Sci.\ \textbf{2470}, Springer, 2002, 763.
\href{https://doi.org/10.1007/3-540-46135-3_61}{doi:10.1007/3-540-46135-3\_61}.
\bibitem{Picard} J.-C.~Picard, \emph{Maximal closure of a graph and applications to combinatorial
problems}, Management Sci.\ \textbf{22} (1976), 1268--1272.
\href{https://doi.org/10.1287/mnsc.22.11.1268}{doi:10.1287/mnsc.22.11.1268}.
\bibitem{PB} S.~Prestwich, J.~C.~Beck, \emph{Exploiting dominance in three symmetric problems}, in:
Proceedings of the Fourth International Workshop on Symmetry and Constraint Satisfaction Problems
(SymCon'04), 2004, 63--70.
\bibitem{SciPy} P.~Virtanen et al., \emph{SciPy 1.0: fundamental algorithms for scientific computing in
Python}, Nature Methods \textbf{17} (2020), 261--272.
\href{https://doi.org/10.1038/s41592-019-0686-2}{doi:10.1038/s41592-019-0686-2}.
\bibitem{SPG} B.~M.~Smith, K.~E.~Petrie, I.~P.~Gent, \emph{Models and symmetry breaking for `peaceable
armies of queens'}, in: CPAIOR 2004, Lecture Notes in Comput.\ Sci.\ \textbf{3011}, Springer, 2004,
271--286. \href{https://doi.org/10.1007/978-3-540-24664-0_19}{doi:10.1007/978-3-540-24664-0\_19}.
\bibitem{vBM} M.~F.~van Bommel, K.~T.~MacEachern, \emph{Armies of chess queens}, Math.\ Intelligencer
\textbf{40}(2) (2018), 10--15. \href{https://doi.org/10.1007/s00283-017-9771-5}{doi:10.1007/s00283-017-9771-5}.
\bibitem{DRAT} N.~Wetzler, M.~J.~H.~Heule, W.~A.~Hunt, \emph{DRAT-trim: efficient checking and trimming
using expressive clausal proofs}, in: SAT 2014, Lecture Notes in Comput.\ Sci.\ \textbf{8561}, Springer,
2014, 422--429. \href{https://doi.org/10.1007/978-3-319-09284-3_31}{doi:10.1007/978-3-319-09284-3\_31}.
\bibitem{YZ} Y.~Yao, D.~Zeilberger, \emph{Numerical and symbolic studies of the peaceable queens
problem}, Exp.\ Math.\ \textbf{31} (2022), 269--279; arXiv:1902.05886.
\href{https://doi.org/10.1080/10586458.2019.1616338}{doi:10.1080/10586458.2019.1616338}.
\end{thebibliography}
\end{document}
